% Section 2.2.1 Architectural Approaches
% Owner: Slice 4 (issue #20).
\subsubsection{Architectural Approaches}
\label{sec:approaches}

The \POSAdapter translates sales for the \StoreAgent, whose local transaction commits
inventory, source sequence and \Outbox event together: Transactions (Bass, Clements,
and Kazman 61). The \Outbox retries committed events with unchanged IDs until
acknowledged: Retry (60). Escalation retains undelivered records.

The \CFE atomically persists each message, its identity and source position before
returning an \texttt{accepted} \Acknowledgement and distributing it to subscribers.
A known retry receives another acknowledgement without a second log entry.
A \StoreAgent raises a
\ShippingOrder when stock reaches its configured minimum; the store manager confirms
it before the \Outbox queues it for delivery. \ShippingOrderFulfillment receives the
confirmed order and presents it to the \DC manager for fulfilment. When dispatch
changes warehouse stock, the \DCAgent publishes an \InventoryEvent through the
\CFE for \CentralMasterInventory to apply.

The Publish-Subscribe pattern (Bass, Clements, and Kazman 129--130) uses the \CFE
as its event bus and realises Use an intermediary (107; 123). The \IEP and
message definitions form a shared local standard: Adhere to standards (107).
The publisher saves
\texttt{acceptedThrough} before pruning its \Outbox. Each subscriber sends cumulative
\texttt{appliedThrough} only after its effects, duplicate records and cursor commit
together; irrelevant messages still advance its durable processing position.

\paragraph{Routing and retention.}
Subscriptions select authorized source streams; consumers filter locally. Dropping
family topics avoids maintaining routing when product-family membership changes,
but increases fan-out, network traffic and client filtering. Consumers acknowledge
their progress even through messages they ignore.
The \CFE prunes a message only after every relevant active subscription acknowledges
through its sequence, then a further 60 minutes from the last required acknowledgement
(Section~\ref{sec:iep}). With no relevant subscription, the window starts at acceptance.
Disconnected subscriptions pin unread history. Slow consumers therefore increase
storage and can delay publishers through backpressure; the timer cannot expire events
still owed to them.

\paragraph{Alternatives and trade-offs.}
Synchronous \DC confirmation would make local updates depend on the WAN.
Local commitment and asynchronous forwarding remove that dependency; central fan-out
avoids a full broker per store. Costs are queues, retries and temporary central
staleness. Delivery assumes recoverable records, sufficient capacity and eventual reconnection.

\paragraph{Recovery decisions.}
The C\&C and Allocation views (Sections~\ref{sec:cc-view} and~\ref{sec:allocation-view})
select supervised \StoreAgent restart and rebuild from its local log/snapshot,
preserved separately from pruned \Outbox transmissions. Failure detection through
restored access must meet the 30-minute local-outage limit.
\CentralMasterInventory restarts from its committed cursor and retained CFE history.
The single \DC host needs replacement/backup recovery for host or storage failure.
\opResync compares durable positions before replay: expired identities require
reconciliation; missing history requires a consistent snapshot/backup.
Recovery never resets stock or reuses source sequences.

\paragraph{Scenario tactics and patterns.}
The tables map every scenario to mechanisms and costs. Tactic page numbers refer to
Bass, Clements, and Kazman, fourth edition. The (H,M) mechanisms are selected in the
named views; capacity and recovery figures remain targets to validate. Other rows
marked \emph{proposed} or \emph{options} identify remaining refinements.

\begingroup
\footnotesize
\setlength{\tabcolsep}{3pt}
\renewcommand{\arraystretch}{1.0}
\noindent\textbf{Named patterns and supporting tactics.}\par
\noindent
\begin{tabularx}{\linewidth}{|>{\raggedright\arraybackslash}p{0.28\linewidth}|>{\raggedright\arraybackslash}p{0.23\linewidth}|>{\raggedright\arraybackslash}X|}
\hline
\textbf{Scenarios} & \textbf{Pattern and source} & \textbf{Supporting tactic / realisation}\\\hline
\scnref{AV-2}, \scnref{PERF-2}, \scnref{SCAL-2}, \scnref{MOD-1}, \scnref{INT-1}, \scnref{INT-2}
& Publish-Subscribe (Bass et al. 129--130)
& Use an intermediary: the \CFE event bus separates publishers and consumers (Section~\ref{sec:cc-view}); fan-out adds central work and dependency.\\\hline
\scnref{MOD-2}, \scnref{PERF-3}
& Composite (Gamma et al. 163)
& Encapsulate: Common's \texttt{CatalogItem}, \texttt{ProductFamily} and \texttt{Product} support recursive value queries (Figure~\ref{fig:domain-model}); traversal must guard cycles and repeated products.\\\hline
\end{tabularx}

\par\medskip\noindent
\begin{tabularx}{\linewidth}{|>{\raggedright\arraybackslash}p{0.09\linewidth}|>{\raggedright\arraybackslash}p{0.22\linewidth}|>{\raggedright\arraybackslash}X|>{\raggedright\arraybackslash}p{0.20\linewidth}|}
\hline
\textbf{ID} & \textbf{Tactic (4e, page)} & \textbf{Mechanism / view} & \textbf{Cost}\\\hline
\tacticrow{\scnref{AV-1}}
  {Graceful degradation (60); Escalating restart (60--61)}
  {Local inventory and \Outbox buffering; supervised \StoreAgent rebuild from its local log (Sections~\ref{sec:cc-view}, \ref{sec:allocation-view}).}
  {Central state lags; storage capacity limits disconnected operation.}
\tacticrow{\scnref{AV-2}}
  {Transactions (61); Retry (60)}
  {Same-ID retries; atomic state/progress; publisher \texttt{accepted} and subscriber cumulative \texttt{applied} acks (Section~\ref{sec:iep}).}
  {Persistence and duplicate tracking; \CFE history waits for all required acks plus 60 minutes.}
\tacticrow{\scnref{AV-3}}
  {Transactions (61)}
  {\CentralMasterInventory commits state/cursor together; restart and \opReplay use history held until all required cumulative acks plus 60 minutes (Sections~\ref{sec:cc-view}, \ref{sec:allocation-view}).}
  {Replay delay; disconnected consumers pin history and can trigger backpressure.}
\end{tabularx}\par\medskip

\noindent
\begin{tabularx}{\linewidth}{|>{\raggedright\arraybackslash}p{0.09\linewidth}|>{\raggedright\arraybackslash}p{0.22\linewidth}|>{\raggedright\arraybackslash}X|>{\raggedright\arraybackslash}p{0.20\linewidth}|}
\hline
\textbf{ID} & \textbf{Tactic (4e, page)} & \textbf{Mechanism / view} & \textbf{Cost}\\\hline
\tacticrow{\scnref{PERF-1}}
  {Introduce concurrency (141)}
  {Local transactions overlap ordered \Outbox transmission (Sections~\ref{sec:cc-view}, \ref{sec:allocation-view}).}
  {Coordination and resource overhead; validate 100 events per second.}
\tacticrow{\scnref{PERF-2}}
  {Introduce concurrency (141); Bound queue sizes (142)}
  {Independent consumer delivery; bounded admission returns \texttt{BUSY} while publishers retain pending events (Section~\ref{sec:iep}).}
  {Contention and backpressure; recovery needs spare capacity.}
\tacticrow{\scnref{PERF-3}}
  {Increase efficiency of resource usage (141)}
  {Proposed indexed inventory queries; Common's recursive Composite traversal (Figure~\ref{fig:domain-model}).}
  {Index maintenance, storage and correct shared-product counting.}
\tacticrow{\scnref{PERF-4}}
  {Introduce concurrency; Increase resources (141)}
  {\Outbox: 180{,}000-event WAN buffer; catch-up capacity 2{,}000 events per second while live sales continue (Sections~\ref{sec:cc-view}, \ref{sec:allocation-view}).}
  {Local storage and added WAN/\DC headroom; replay competes with live work.}
\end{tabularx}\par\medskip

\noindent
\begin{tabularx}{\linewidth}{|>{\raggedright\arraybackslash}p{0.09\linewidth}|>{\raggedright\arraybackslash}p{0.22\linewidth}|>{\raggedright\arraybackslash}X|>{\raggedright\arraybackslash}p{0.20\linewidth}|}
\hline
\textbf{ID} & \textbf{Tactic (4e, page)} & \textbf{Mechanism / view} & \textbf{Cost}\\\hline
\tacticrow{\scnref{SEC-1}}
  {Authenticate actors; Authorize actors (174)}
  {Options for manager/agent identities, inventory access checks and restricted database accounts.}
  {Identity infrastructure and permission administration.}
\tacticrow{\scnref{SEC-2}}
  {Encrypt data (175); Verify message integrity (174); Authenticate actors (174)}
  {Store--\DC \IEP uses TLS 1.3 and mutual certificates with provisioned trust (Sections~\ref{sec:cc-view}, \ref{sec:allocation-view}).}
  {Certificate provisioning/revocation and cryptographic overhead.}
\tacticrow{\scnref{SEC-3}}
  {Authorize actors (174)}
  {Proposed operation and inventory-scope checks in \StoreAgent and \CentralMasterInventory access paths.}
  {Accurate scope assignments and consistent enforcement.}
\tacticrow{\scnref{SCAL-1}}
  {Defer binding (configuration time, 124)}
  {One configured \StoreAgent, registered source identity and certificates per new store (Sections~\ref{sec:cc-view}, \ref{sec:allocation-view}).}
  {Provisioning work; central throughput needs separate validation.}
\tacticrow{\scnref{SCAL-2}}
  {Use an intermediary (107)}
  {\CFE registers conforming subscribers through \opSubscribe at runtime.}
  {Each client adds distribution work and possible backlog.}
\end{tabularx}\par\medskip

\noindent
\begin{tabularx}{\linewidth}{|>{\raggedright\arraybackslash}p{0.09\linewidth}|>{\raggedright\arraybackslash}p{0.22\linewidth}|>{\raggedright\arraybackslash}X|>{\raggedright\arraybackslash}p{0.20\linewidth}|}
\hline
\textbf{ID} & \textbf{Tactic (4e, page)} & \textbf{Mechanism / view} & \textbf{Cost}\\\hline
\tacticrow{\scnref{USE-1}}
  {Maintain task model (201)}
  {Order context and confirmation through \StoreAgent before \Outbox publication.}
  {Interaction-state logic; validate the minimal-training criterion.}
\tacticrow{\scnref{USE-2}}
  {Maintain task model (201)}
  {Proposed \ReportingClient guidance based on report type, inputs and scope.}
  {Per-report interaction logic and user evaluation.}
\end{tabularx}\par\medskip

\noindent
\begin{tabularx}{\linewidth}{|>{\raggedright\arraybackslash}p{0.09\linewidth}|>{\raggedright\arraybackslash}p{0.22\linewidth}|>{\raggedright\arraybackslash}X|>{\raggedright\arraybackslash}p{0.20\linewidth}|}
\hline
\textbf{ID} & \textbf{Tactic (4e, page)} & \textbf{Mechanism / view} & \textbf{Cost}\\\hline
\tacticrow{\scnref{MOD-1}}
  {Use an intermediary (107; 123)}
  {New client uses \IEP without changing existing publishers.}
  {Shared-contract compatibility and extra subscriber load.}
\tacticrow{\scnref{MOD-2}}
  {Encapsulate (106; 123)}
  {Common's Composite exposes \texttt{CatalogItem.totalValue()} for leaves and families (Figure~\ref{fig:domain-model}).}
  {Traversal, cycle checks and shared-product counting.}
\tacticrow{\scnref{INT-1}}
  {Use an intermediary; Adhere to standards (107)}
  {The future \AnalyticsClient receives authorized events over the shared \IEP.}
  {Contract compatibility, client provisioning and additional fan-out work.}
\tacticrow{\scnref{INT-2}}
  {Use an intermediary (107)}
  {\CFE subscription interface serves a variable number of conforming clients.}
  {Central dependency and fan-out resource costs.}
\end{tabularx}\par\medskip

\noindent
\begin{tabularx}{\linewidth}{|>{\raggedright\arraybackslash}p{0.09\linewidth}|>{\raggedright\arraybackslash}p{0.22\linewidth}|>{\raggedright\arraybackslash}X|>{\raggedright\arraybackslash}p{0.20\linewidth}|}
\hline
\textbf{ID} & \textbf{Tactic (4e, page)} & \textbf{Mechanism / view} & \textbf{Cost}\\\hline
\tacticrow{\scnref{TEST-1}}
  {Specialized interfaces (188); Sandbox (189)}
  {Proposed isolated test instance with controlled message/acknowledgement loss and observable outcomes.}
  {Test infrastructure; finite cases support only stated assumptions.}
\tacticrow{\scnref{TEST-2}}
  {Specialized interfaces (188); Sandbox (189)}
  {Proposed isolated test instance with failure injection and timing to usable inventory access.}
  {Instrumentation overhead; real recovery time must be measured.}
\end{tabularx}\par\medskip
\endgroup

\paragraph{Validation and remaining risks.}
The selected recovery, buffering, mutual TLS and store-provisioning mechanisms
still need timed recovery and load tests. Full \Outbox admission rejects an update
atomically; central admission returns \texttt{BUSY} before accepting excess work
(Section~\ref{sec:iep}). Credential/scope administration, backup restoration and the
physical dispatch handoff require validation. The \scnref{TEST-2} wording must describe
the restart/recovery mechanism supplied by the deployment.
